\documentclass[10pt,a4paper]{amsart}
\usepackage{amsmath,amssymb,amsthm,graphicx}
% For dealing with references we use the comment environment
\usepackage{verbatim}
\newenvironment{reference}{\comment}{\endcomment}
%\newenvironment{reference}{}{}
\newenvironment{slogan}{\comment}{\endcomment}
\newenvironment{history}{\comment}{\endcomment}
% For commutative diagrams we use Xy-pic
\usepackage[all]{xy}

% We use 2cell for 2-commutative diagrams.
\xyoption{2cell}
\UseAllTwocells
\theoremstyle{plain}
\newtheorem{theorem}[subsection]{Theorem}
\newtheorem{proposition}[subsection]{Proposition}
\newtheorem{lemma}[subsection]{Lemma}

\theoremstyle{definition}
\newtheorem{definition}[subsection]{Definition}
\newtheorem{example}[subsection]{Example}
\newtheorem{exercise}[subsection]{Exercise}
\newtheorem{situation}[subsection]{Situation}

\theoremstyle{remark}
\newtheorem{remark}[subsection]{Remark}
\newtheorem{remarks}[subsection]{Remarks}

\numberwithin{equation}{subsection}
% Macros
%
\def\lim{\mathop{\mathrm{lim}}\nolimits}
\def\colim{\mathop{\mathrm{colim}}\nolimits}
\def\Spec{\mathop{\mathrm{Spec}}}
\def\Hom{\mathop{\mathrm{Hom}}\nolimits}
\def\Ext{\mathop{\mathrm{Ext}}\nolimits}
\def\SheafHom{\mathop{\mathcal{H}\!\mathit{om}}\nolimits}
\def\SheafExt{\mathop{\mathcal{E}\!\mathit{xt}}\nolimits}
\def\Sch{\mathit{Sch}}
\def\Mor{\mathop{\mathrm{Mor}}\nolimits}
\def\Ob{\mathop{\mathrm{Ob}}\nolimits}
\def\Sh{\mathop{\mathit{Sh}}\nolimits}
\def\NL{\mathop{N\!L}\nolimits}
\def\CH{\mathop{\mathrm{CH}}\nolimits}
\def\proetale{{pro\text{-}\acute{e}tale}}
\def\etale{{\acute{e}tale}}
\def\QCoh{\mathit{QCoh}}
\def\Ker{\mathop{\mathrm{Ker}}}
\def\Im{\mathop{\mathrm{Im}}}
\def\Coker{\mathop{\mathrm{Coker}}}
\def\Coim{\mathop{\mathrm{Coim}}}

% Boxtimes
%
\DeclareMathSymbol{\boxtimes}{\mathbin}{AMSa}{"02}

%
% Macros for moduli stacks/spaces
%
\def\QCohstack{\mathcal{QC}\!\mathit{oh}}
\def\Cohstack{\mathcal{C}\!\mathit{oh}}
\def\Spacesstack{\mathcal{S}\!\mathit{paces}}
\def\Quotfunctor{\mathrm{Quot}}
\def\Hilbfunctor{\mathrm{Hilb}}
\def\Curvesstack{\mathcal{C}\!\mathit{urves}}
\def\Polarizedstack{\mathcal{P}\!\mathit{olarized}}
\def\Complexesstack{\mathcal{C}\!\mathit{omplexes}}
% \Pic is the operator that assigns to X its picard group, usage \Pic(X)
% \Picardstack_{X/B} denotes the Picard stack of X over B
% \Picardfunctor_{X/B} denotes the Picard functor of X over B
\def\Pic{\mathop{\mathrm{Pic}}\nolimits}
\def\Picardstack{\mathcal{P}\!\mathit{ic}}
\def\Picardfunctor{\mathrm{Pic}}
\def\Deformationcategory{\mathcal{D}\!\mathit{ef}}
\usepackage{cleveref}
\crefname{originalsection}{original section}{original sections}
\Crefname{originalsection}{Original section}{Original sections}
\makeatletter
\newlabel{section-notation}{{043A}{source}{}{sourcepointer-1}{}}
\newlabel{section-notation@cref}{{[originalsection][0][]043A}{[1][0][]source}{}{}{}}
\makeatother
\crefname{originallemma}{original lemma}{original lemmas}
\Crefname{originallemma}{Original lemma}{Original lemmas}
\makeatletter
\newlabel{groupoids-lemma-colimit-kappa}{{077U}{source}{}{sourcepointer-2}{}}
\newlabel{groupoids-lemma-colimit-kappa@cref}{{[originallemma][0][]077U}{[1][0][]source}{}{}{}}
\makeatother
\crefname{originalproposition}{original proposition}{original propositions}
\Crefname{originalproposition}{Original proposition}{Original propositions}
\makeatletter
\newlabel{spaces-properties-proposition-quasi-coherent}{{03M3}{source}{}{sourcepointer-3}{}}
\newlabel{spaces-properties-proposition-quasi-coherent@cref}{{[originalproposition][0][]03M3}{[1][0][]source}{}{}{}}
\makeatother
\crefname{originallemma}{original lemma}{original lemmas}
\Crefname{originallemma}{Original lemma}{Original lemmas}
\makeatletter
\newlabel{lemma-crystals-in-quasi-coherent-modules}{{077Z}{source}{}{sourcepointer-4}{}}
\newlabel{lemma-crystals-in-quasi-coherent-modules@cref}{{[originallemma][0][]077Z}{[1][0][]source}{}{}{}}
\makeatother
\crefname{originallemma}{original lemma}{original lemmas}
\Crefname{originallemma}{Original lemma}{Original lemmas}
\makeatletter
\newlabel{properties-lemma-set-of-iso-classes}{{077M}{source}{}{sourcepointer-5}{}}
\newlabel{properties-lemma-set-of-iso-classes@cref}{{[originallemma][0][]077M}{[1][0][]source}{}{}{}}
\makeatother
\title{Lemma set generators in spaces-groupoids.tex}
\author{The Stacks Project Authors}
\date{stacks-project:a04446e57ec1fbc252a871afcec7752fb2807b14; GNU Free Documentation License 1.2 (as published in the Stacks Project COPYING)}
\begin{document}
\maketitle
\noindent\textit{Editorial note.} All mathematical statement and proof text is reproduced verbatim from the author's chapter spaces-groupoids.tex of the Stacks Project, fixed at commit a04446e57ec1fbc252a871afcec7752fb2807b14. The pinned preamble is supplied as its own file; the wrapper class, the xr-hyper wiring and the \textbackslash{}externaldocument list are omitted because the assembler supplies the class.\par\medskip
\noindent\textit{Editorial note.} Numbering inside this standalone excerpt is editorial, not the full Stacks book print numbering.\par\medskip
\noindent\textit{Editorial note.} The candidate is from the corrected band-I follow-on screen (line C) of the extractor's stacks-candidates line: its author label is absent from the live master in every form searched and it belongs to none of the earlier judged pools. The H2 judgement and its reason were made in this module by reading the author's own statement and complete proof in the fixed checkout. Original candidate key: stacks:spaces-groupoids.tex:lemma-set-generators.\par\medskip
\noindent\textit{Editorial note.} Reference labels in the proof resolve to navigation pointers to the original author source; these pointers do not reproduce proofs of the prerequisites.\par\medskip
\noindent\textit{Allowed standard prerequisites (editorial).}\begin{itemize}
\item Prerequisites named by the author inside the counted span. The counted span is reproduced verbatim, including the author's own citations. Its prerequisites are the earlier results of the same fixed Stacks Project version that the author cites there and in the chapter's lead-in.
\end{itemize}\medskip
\ifdefined\hypertarget\hypertarget{sourcepointer-1}{}\fi
\noindent\textit{Source reference.} Original section 043A: sources/spaces-groupoids.tex, source lines 55--55. Navigation pointer used in the counted span; the span contains the pointer only and the cited result is not reproduced inside it. Registered as a non-load-bearing navigation stub in the established form for Stacks material. No bibliography citation is involved.\par\medskip
\ifdefined\hypertarget\hypertarget{sourcepointer-2}{}\fi
\noindent\textit{Source reference.} Original lemma 077U: sources/groupoids.tex, source lines 3153--3153. Navigation pointer used in the counted span; the span contains the pointer only and the cited result is not reproduced inside it. Registered as a non-load-bearing navigation stub in the established form for Stacks material. No bibliography citation is involved.\par\medskip
\ifdefined\hypertarget\hypertarget{sourcepointer-3}{}\fi
\noindent\textit{Source reference.} Original proposition 03M3: sources/spaces-properties.tex, source lines 5727--5727. Navigation pointer used in the counted span; the span contains the pointer only and the cited result is not reproduced inside it. Registered as a non-load-bearing navigation stub in the established form for Stacks material. No bibliography citation is involved.\par\medskip
\ifdefined\hypertarget\hypertarget{sourcepointer-4}{}\fi
\noindent\textit{Source reference.} Original lemma 077Z: sources/spaces-groupoids.tex, source lines 1496--1496. Navigation pointer used in the counted span; the span contains the pointer only and the cited result is not reproduced inside it. Registered as a non-load-bearing navigation stub in the established form for Stacks material. No bibliography citation is involved.\par\medskip
\ifdefined\hypertarget\hypertarget{sourcepointer-5}{}\fi
\noindent\textit{Source reference.} Original lemma 077M: sources/properties.tex, source lines 3532--3532. Navigation pointer used in the counted span; the span contains the pointer only and the cited result is not reproduced inside it. Registered as a non-load-bearing navigation stub in the established form for Stacks material. No bibliography citation is involved.\par\medskip
\begin{lemma}
\label{lemma-set-generators}
Let $B \to S$ as in Section \ref{section-notation}.
Let $(U, R, s, t, c)$ be a groupoid in algebraic spaces over $B$.
If $s$, $t$ are flat, then there exists a set $T$ and a family of objects
$(\mathcal{F}_t, \alpha_t)_{t \in T}$ of $\QCoh(U, R, s, t, c)$
such that every object $(\mathcal{F}, \alpha)$ is the directed colimit
of its submodules isomorphic to one of the objects $(\mathcal{F}_t, \alpha_t)$.
\end{lemma}

\begin{proof}
This lemma is a generalization of
Groupoids, Lemma \ref{groupoids-lemma-colimit-kappa}
which deals with the case of a groupoid in schemes.
We can't quite use the same argument, so we use the
material on ``crystals of quasi-coherent sheaves'' we developed above.

\medskip\noindent
Choose a scheme $W$ and a surjective \'etale morphism $W \to U$.
Choose a scheme $V$ and a surjective \'etale morphism
$V \to W \times_{U, s} R$. Choose a scheme $V'$
and a surjective \'etale morphism $V' \to R \times_{t, U} W$.
Consider the collection of schemes
$$
I = \{W, W \times_U W, V, V', V \times_R V'\}
$$
and the set of morphisms of schemes
$$
\Phi = \{\text{pr}_i : W \times_U W \to W, V \to W, V' \to W,
V \times_R V' \to V, V \times_R V' \to V'\}
$$
Set $X = (I, \Phi)$. Recall that we have defined a category $\textit{CQC}(X)$
of crystals of quasi-coherent sheaves on $X$. There is a functor
$$
\QCoh(U, R, s, t, c) \longrightarrow \textit{CQC}(X)
$$
which assigns to $(\mathcal{F}, \alpha)$ the sheaf $\mathcal{F}|_W$ on $W$,
the sheaf $\mathcal{F}|_{W \times_U W}$ on $W \times_U W$, the
pullback of $\mathcal{F}$ via $V \to W \times_{U, s} R \to W \to U$
on $V$, the pullback of $\mathcal{F}$ via
$V' \to R \times_{t, U} W \to W \to U$ on $V'$, and finally the
pullback of $\mathcal{F}$ via
$V \times_R V' \to V \to W \times_{U, s} R \to W \to U$ on $V \times_R V'$.
As comparison maps $\{\alpha_\phi\}_{\phi \in \Phi}$
we use the obvious ones (coming from associativity of pullbacks)
except for the map
$\phi = \text{pr}_{V'} : V \times_R V' \to V'$ we use the pullback
of $\alpha : t^*\mathcal{F} \to s^*\mathcal{F}$ to $V \times_R V'$.
This makes sense because of the following commutative diagram
$$
\xymatrix{
& V \times_R V' \ar[ld] \ar[rd] \\
V \ar[rd] \ar[dd] & & V' \ar[ld] \ar[dd] \\
& R \ar@<-1ex>[dd]_s \ar@<1ex>[dd]^t \\
W \ar[rd] & & W \ar[ld] \\
& U
}
$$
The functor displayed above isn't an equivalence of categories.
However, since $W \to U$ is surjective \'etale it is faithful\footnote{In
fact the functor is fully faithful, but we won't need this.}.
Since all the morphisms in the diagram above are flat we see that
it is an exact functor of abelian categories. Moreover, we
claim that given $(\mathcal{F}, \alpha)$ with image
$(\{\mathcal{F}_i\}_{i \in I}, \{\alpha_\phi\}_{\phi \in \Phi})$
there is a $1$-to-$1$ correspondence between quasi-coherent submodules
of $(\mathcal{F}, \alpha)$ and
$(\{\mathcal{F}_i\}_{i \in I}, \{\alpha_\phi\}_{\phi \in \Phi})$.
Namely, given a submodule of
$(\{\mathcal{F}_i\}_{i \in I}, \{\alpha_\phi\}_{\phi \in \Phi})$
compatibility of the submodule over $W$ with the projection maps
$W \times_U W \to W$ will guarantee the submodule comes from a
quasi-coherent submodule of $\mathcal{F}$ (by
Properties of Spaces, Proposition
\ref{spaces-properties-proposition-quasi-coherent})
and compatibility with $\alpha_{\text{pr}_{V'}}$ will insure this
subsheaf is compatible with $\alpha$ (details omitted).

\medskip\noindent
Choose a cardinal $\kappa$ as in
Lemma \ref{lemma-crystals-in-quasi-coherent-modules}
for the system $X = (I, \Phi)$. It is clear from
Properties, Lemma \ref{properties-lemma-set-of-iso-classes}
that there is a set of isomorphism classes of $\kappa$-generated
crystals in quasi-coherent sheaves on $X$.
Hence the result is clear.
\end{proof}

\end{document}
